% =====================================================================
% pertemuan-09.tex
% Pertemuan 9 - Top-Down Parsing: Recursive Descent, LL(1), FIRST, FOLLOW
% Sub-CPMK 3.3, 4.1 - Praktikum 2
% =====================================================================
\documentclass[aspectratio=169,10pt,t]{beamer}
\usepackage{config/beamerthemeubb}

\UBBmeeting{9}

\begin{document}

\UBBtitleframe{9}{Top-Down Parsing: Recursive Descent, LL(1), FIRST, dan FOLLOW}{Praktikum 2}

% ---------------------------------------------------------------------
\begin{frame}{Capaian Pertemuan}
  \textbf{CPMK induk:} CPMK 03 --- Merancang tata bahasa bebas konteks
  untuk bahasa sederhana; CPMK 04 --- Mengimplementasikan parser
  sederhana untuk validasi kode.

  \vspace{2pt}
  \begin{block}{Sub-CPMK yang dinilai}
    \footnotesize
    \begin{itemize}
      \setlength{\itemsep}{1pt}
      \item \textbf{Sub-CPMK 3.3:} Membuat tabel himpunan FIRST dan
            FOLLOW untuk grammar subset bahasa C dan menguji keputusan
            parser LL(1) berdasarkan tabel tersebut.
      \item \textbf{Sub-CPMK 4.1:} Mengimplementasikan
            \textit{recursive descent parser} dan \textit{predictive
            parser} LL(1) dalam bahasa C, dengan berkas kode dan
            keluaran validasi sebagai bukti.
    \end{itemize}
  \end{block}
\end{frame}

% ---------------------------------------------------------------------
\begin{frame}{Agenda Pertemuan (150 menit)}
  \small
  \begin{tabular}{@{}p{1.4cm}p{12.3cm}@{}}
    \toprule
    \textbf{Menit} & \textbf{Kegiatan} \\
    \midrule
    20 & Briefing Praktikum 2 dan cakupan artefak \\
    25 & Ceramah: recursive descent dan predictive parsing \\
    30 & Latihan: menghitung FIRST dan FOLLOW \\
    30 & Lokakarya: tabel LL(1) dan penelusuran stack \\
    35 & Praktik: parser C, AST, dan diagnostik \\
    10 & Pengujian, laporan, dan penutup \\
    \bottomrule
  \end{tabular}

  \vspace{6pt}
  {\scriptsize\color{ubbblue!75}Porsi praktik dan latihan 105 menit;
  ceramah 25 menit. Setiap mahasiswa menyerahkan tabel dan parser.}
\end{frame}

% ---------------------------------------------------------------------
\begin{frame}{Posisi Pertemuan pada Rantai Proyek}
  \UBBchain{3}

  \vspace{2pt}
  \begin{itemize}
    \item Pertemuan ini membangun simpul \textbf{parser}: dari grammar
          bersih menjadi pengenal struktur program.
    \item Keluaran parser berupa AST menjadi masukan analisis semantik
          pada Pertemuan 13--15.
    \item Praktikum 2 menilai Sub-CPMK 3.1--3.4 dan 4.1.
  \end{itemize}
\end{frame}

% ---------------------------------------------------------------------
\begin{frame}{Recursive Descent: Satu Fungsi per Nonterminal}
  \begin{columns}[T,onlytextwidth]
    \begin{column}{0.58\textwidth}
      \begin{itemize}
        \item Satu fungsi untuk setiap nonterminal.
        \item \texttt{parseE} mengenali string dari \(E\);
              \texttt{parseT} dari \(T\).
        \item Setiap fungsi memilih aturan berdasarkan token saat ini.
        \item Mudah ditulis, mudah di-\textit{debug}, dan pesan kesalahan
              baik.
        \item Syarat: tanpa \textit{left recursion}, sudah difaktorkan,
              satu token lookahead.
      \end{itemize}
    \end{column}
    \begin{column}{0.38\textwidth}
      \centering
      \begin{tikzpicture}[
          f/.style={draw=ubbblue!40,fill=ubbblue!7,rounded corners=3pt,
            minimum width=3.4cm,minimum height=0.52cm,font=\scriptsize,
            text=ubbblue},
          ar/.style={-{Stealth[length=1.4mm]},draw=ubbblue!55}
        ]
        \node[f] (p) {\texttt{parser\_program}};
        \node[f,below=0.18cm of p] (t) {\texttt{parse\_topdecl}};
        \node[f,below=0.18cm of t] (b) {\texttt{parse\_block}};
        \node[f,below=0.18cm of b] (s) {\texttt{parse\_stmt}};
        \node[f,below=0.18cm of s] (e) {\texttt{parse\_expr}};
        \node[f,below=0.18cm of e] (pr) {\texttt{parse\_primary}};
        \foreach \a/\b in {p/t,t/b,b/s,s/e,e/pr}{\draw[ar] (\a)--(\b);}
      \end{tikzpicture}
    \end{column}
  \end{columns}
\end{frame}

% ---------------------------------------------------------------------
\begin{frame}{Struktur Parser dan Primitifnya}
  \begin{block}{Struktur \texttt{Parser}}
    \small
    Menyimpan \texttt{Lexer}, token saat ini (\texttt{current}), dan
    status kesalahan (\texttt{had\_error}).
  \end{block}

  \begin{itemize}
    \item \texttt{advance()}: mengambil token berikutnya dari lexer.
    \item \texttt{check(type)}: menguji jenis token tanpa
          mengonsumsinya.
    \item \texttt{match(type)}: mengonsumsi token jika jenisnya sesuai.
    \item Ketiganya menjadi primitif seluruh parser.
  \end{itemize}

  \vspace{3pt}
  {\footnotesize\color{ubbblue!80}Jika token tidak cocok, parser
  melaporkan kesalahan dan mengembalikan simpul pengganti agar
  analisis dapat dilanjutkan.}
\end{frame}

% ---------------------------------------------------------------------
\begin{frame}[fragile]{Contoh Fungsi: \texttt{parse\_primary}}
  \begin{columns}[T,onlytextwidth]
    \begin{column}{0.58\textwidth}
      \lstinputlisting[language=C,firstline=107,lastline=121,
        firstnumber=107,
        caption={Awal \texttt{parse\_primary} pada \file{parser.c}},
        label={lst:p9-primary}]{../buku_ajar/code/subsetc/src/parser.c}
    \end{column}
    \begin{column}{0.38\textwidth}
      \begin{itemize}
        \item Literal integer, float, char, dan string.
        \item Identifier, pemanggilan, dan ekspresi dalam kurung
              ditangani cabang berikutnya.
        \item Setiap cabang memeriksa token saat ini.
        \item Tidak ada cabang cocok: kesalahan sintaksis.
      \end{itemize}
    \end{column}
  \end{columns}
\end{frame}

% ---------------------------------------------------------------------
\begin{frame}{Himpunan FIRST}
  \begin{block}{Definisi}
    \small
    \(\text{FIRST}(\alpha)\) adalah himpunan terminal yang dapat muncul
    sebagai simbol pertama pada string turunan \(\alpha\). Jika
    \(\alpha\) dapat menurunkan \(\varepsilon\), maka
    \(\varepsilon \in \text{FIRST}(\alpha)\).
  \end{block}

  \begin{enumerate}
    \setlength{\itemsep}{1pt}
    \item Jika \(X\) terminal, \(\text{FIRST}(X)=\{X\}\).
    \item Jika \(X \to \varepsilon\), tambahkan \(\varepsilon\).
    \item Jika \(X \to Y_1Y_2\cdots Y_k\), tambahkan
          \(\text{FIRST}(Y_1)\setminus\{\varepsilon\}\); jika \(Y_1\)
          nullable, lanjut ke \(Y_2\), dan seterusnya.
    \item Ulangi sampai tidak ada perubahan.
  \end{enumerate}
\end{frame}

% ---------------------------------------------------------------------
\begin{frame}{Himpunan FOLLOW}
  \begin{block}{Definisi}
    \small
    \(\text{FOLLOW}(A)\) adalah himpunan terminal yang dapat muncul
    tepat setelah \(A\) pada suatu bentuk kalimat. Simbol akhir
    masukan \(\$\) ditambahkan ke FOLLOW simbol awal.
  \end{block}

  \begin{enumerate}
    \setlength{\itemsep}{1pt}
    \item Tambahkan \(\$\) ke \(\text{FOLLOW}(S)\).
    \item Untuk \(A \to \alpha B\beta\), tambahkan
          \(\text{FIRST}(\beta)\setminus\{\varepsilon\}\) ke
          \(\text{FOLLOW}(B)\).
    \item Jika \(\beta\) nullable atau kosong, tambahkan
          \(\text{FOLLOW}(A)\) ke \(\text{FOLLOW}(B)\).
    \item Ulangi sampai stabil.
  \end{enumerate}
\end{frame}

% ---------------------------------------------------------------------
\begin{frame}{Tabel FIRST dan FOLLOW (Artefak Sub-CPMK 3.3)}
  \small
  \setlength{\tabcolsep}{6pt}
  \begin{tabular}{@{}l l l@{}}
    \toprule
    \textbf{Nonterminal} & \textbf{FIRST} & \textbf{FOLLOW} \\
    \midrule
    \(E\)  & \(\{\code{id}, (\}\) & \(\{\$, )\}\) \\
    \(E'\) & \(\{+, \varepsilon\}\) & \(\{\$, )\}\) \\
    \(T\)  & \(\{\code{id}, (\}\) & \(\{\$, ), +\}\) \\
    \(T'\) & \(\{*, \varepsilon\}\) & \(\{\$, ), +\}\) \\
    \(F\)  & \(\{\code{id}, (\}\) & \(\{\$, ), *, +\}\) \\
    \bottomrule
  \end{tabular}

  \vspace{6pt}
  \begin{block}{Grammar yang dihitung}
    \footnotesize
    \(E \to TE'\), \(E' \to +TE' \mid \varepsilon\),
    \(T \to FT'\), \(T' \to *FT' \mid \varepsilon\),
    \(F \to (E) \mid \code{id}\).
  \end{block}

  \vspace{2pt}
  {\footnotesize\color{ubbblue!80}Setiap nilai dihitung dari aturan
  grammar, bukan dari intuisi.}
\end{frame}

% ---------------------------------------------------------------------
\begin{frame}{Keputusan LL(1)}
  \begin{block}{Syarat LL(1)}
    \small
    Untuk setiap nonterminal \(A\) dan pasangan aturan
    \(A \to \alpha\), \(A \to \beta\):
    \(\text{FIRST}(\alpha) \cap \text{FIRST}(\beta)=\varnothing\); jika
    salah satu nullable,
    \(\text{FIRST}(\alpha) \cap \text{FOLLOW}(A)=\varnothing\).
  \end{block}

  \small
  \setlength{\tabcolsep}{4pt}
  \begin{tabular}{@{}p{5.4cm}p{2.4cm}p{5.4cm}@{}}
    \toprule
    \textbf{Keadaan grammar} & \textbf{Sel ganda} & \textbf{Keputusan} \\
    \midrule
    Grammar ekspresi ini & tidak ada & LL(1) \\
    Grammar dengan \textit{left recursion} & ada &
      bukan LL(1) sebelum transformasi \\
    Grammar dengan prefiks umum & ada &
      bukan LL(1) sebelum factoring \\
    \bottomrule
  \end{tabular}
\end{frame}

% ---------------------------------------------------------------------
\begin{frame}{Tabel Parsing LL(1)}
  \begin{block}{Konstruksi}
    \footnotesize
    Untuk aturan \(A \to \alpha\) dan \(a \in \text{FIRST}(\alpha)\),
    isi \(M[A,a]=A \to \alpha\). Jika
    \(\varepsilon \in \text{FIRST}(\alpha)\), isi
    \(M[A,b]=A \to \alpha\) untuk setiap \(b \in \text{FOLLOW}(A)\).
    Sel kosong berarti kesalahan sintaksis.
  \end{block}

  \scriptsize
  \begin{tabular}{@{}l l l l l@{}}
    \toprule
     & \texttt{id} & \texttt{+} & \texttt{*} & \texttt{\$} \\
    \midrule
    \(E\)  & \(E \to TE'\) & & & \\
    \(E'\) & & \(E' \to +TE'\) & & \(E' \to \varepsilon\) \\
    \(T\)  & \(T \to FT'\) & & & \\
    \(T'\) & & \(T' \to \varepsilon\) & \(T' \to *FT'\) &
      \(T' \to \varepsilon\) \\
    \(F\)  & \(F \to \code{id}\) & & & \\
    \bottomrule
  \end{tabular}

  \vspace{4pt}
  {\scriptsize Kolom \texttt{(} dan \texttt{)} diisi aturan yang sama
  untuk \(E,T\) dan \(\varepsilon\) untuk \(E',T'\).}
\end{frame}

% ---------------------------------------------------------------------
\begin{frame}{Algoritma Predictive Parser}
  \begin{enumerate}
    \setlength{\itemsep}{1pt}
    \item Dorong \(\$\) dan simbol awal ke stack.
    \item \(a\) adalah token saat ini.
    \item Selama stack tidak kosong:
          \begin{enumerate}
            \item \(X\) adalah puncak stack;
            \item jika \(X\) terminal atau \(\$\), cocokkan dengan
                  \(a\); jika cocok, pop dan majukan \(a\);
            \item jika \(X\) nonterminal, lihat \(M[X,a]\);
            \item jika aturan \(X \to Y_1\cdots Y_k\) ada, pop \(X\)
                  dan dorong \(Y_k,\ldots,Y_1\);
            \item jika sel kosong, laporkan kesalahan.
          \end{enumerate}
  \end{enumerate}

  \vspace{2pt}
  {\footnotesize\color{ubbblue!80}Parser tidak pernah mundur; waktu
  linear terhadap panjang masukan.}
\end{frame}

% ---------------------------------------------------------------------
\begin{frame}{Penelusuran Stack untuk \texttt{id + id * id}}
  \scriptsize
  \setlength{\tabcolsep}{4pt}
  \begin{tabular}{@{}p{2.6cm}p{4.6cm}p{5.8cm}@{}}
    \toprule
    \textbf{Stack} & \textbf{Masukan} & \textbf{Aksi} \\
    \midrule
    \(\$E\) & \(\code{id}+\code{id}*\code{id}\ \$\) & \(E \to TE'\) \\
    \(\$E'T\) & \(\code{id}+\code{id}*\code{id}\ \$\) & \(T \to FT'\) \\
    \(\$E'T'F\) & \(\code{id}+\code{id}*\code{id}\ \$\) & \(F \to \code{id}\) \\
    \(\$E'T'\code{id}\) & \(\code{id}+\code{id}*\code{id}\ \$\) &
      cocokkan \code{id} \\
    \(\$E'T'\) & \(+\code{id}*\code{id}\ \$\) & \(T' \to \varepsilon\) \\
    \(\$E'\) & \(+\code{id}*\code{id}\ \$\) & \(E' \to +TE'\) \\
    \(\$E'T+\) & \(+\code{id}*\code{id}\ \$\) & cocokkan \(+\) \\
    \(\$E'T\) & \(\code{id}*\code{id}\ \$\) & \(T \to FT'\) \\
    \multicolumn{3}{@{}l}{\scriptsize\color{ubbblue!70}\ldots{} penelusuran
      berlanjut sampai stack hanya berisi \(\$\) dan masukan \(\$\)} \\
    \(\$\) & \(\$\) & terima \\
    \bottomrule
  \end{tabular}
\end{frame}

% ---------------------------------------------------------------------
\begin{frame}{Dari Tabel ke Fungsi Parser}
  \begin{itemize}
    \item Aturan \(A \to X_1\cdots X_k\): panggil
          \texttt{parseX1}, \texttt{parseX2}, dan seterusnya.
    \item Terminal: panggil \texttt{match}.
    \item \(\varepsilon\): tidak ada aksi.
    \item Sel kosong: laporkan kesalahan.
  \end{itemize}

  \vspace{4pt}
  \small
  \setlength{\tabcolsep}{4pt}
  \begin{tabular}{@{}p{3.4cm}p{4.4cm}p{5.6cm}@{}}
    \toprule
    \textbf{Nonterminal} & \textbf{Fungsi} & \textbf{Tanggung jawab} \\
    \midrule
    \texttt{program} & \texttt{parser\_program} & seluruh berkas \\
    \texttt{topdecl} & \texttt{parse\_topdecl} & deklarasi fungsi \\
    \texttt{block} & \texttt{parse\_block} & daftar pernyataan \\
    \texttt{stmt} & \texttt{parse\_stmt} & satu pernyataan \\
    \texttt{expr}--\texttt{primary} & \texttt{parse\_expr}, \ldots &
      ekspresi berprioritas \\
    \bottomrule
  \end{tabular}
\end{frame}

% ---------------------------------------------------------------------
\begin{frame}[fragile]{Pemilihan Cabang: \texttt{parse\_stmt}}
  \begin{columns}[T,onlytextwidth]
    \begin{column}{0.58\textwidth}
      \lstinputlisting[language=C,firstline=296,lastline=309,
        firstnumber=296,
        caption={Bagian awal \texttt{parse\_stmt} pada \file{parser.c}},
        label={lst:p9-stmt}]{../buku_ajar/code/subsetc/src/parser.c}
    \end{column}
    \begin{column}{0.38\textwidth}
      \begin{itemize}
        \item Tipe memulai deklarasi lokal.
        \item \texttt{\{} memulai blok.
        \item \texttt{if} memilih \texttt{parse\_if}.
        \item Cabang lain ditangani dengan pola yang sama.
        \item Tidak ada cabang cocok: kesalahan dan simpul pengganti.
      \end{itemize}
    \end{column}
  \end{columns}
\end{frame}

% ---------------------------------------------------------------------
\begin{frame}[fragile]{Sinkronisasi Setelah Kesalahan}
  \begin{columns}[T,onlytextwidth]
    \begin{column}{0.58\textwidth}
      \lstinputlisting[language=C,firstline=42,lastline=55,
        firstnumber=42,
        caption={Fungsi \texttt{synchronize} pada \file{parser.c}},
        label={lst:p9-sync}]{../buku_ajar/code/subsetc/src/parser.c}
    \end{column}
    \begin{column}{0.38\textwidth}
      \begin{itemize}
        \item Majukan token sampai titik aman.
        \item Titik aman: \texttt{;}, \texttt{\}}, atau awal
              pernyataan berikutnya.
        \item Terlalu agresif: laporan palsu.
        \item Terlalu konservatif: kesalahan terlewat.
      \end{itemize}
    \end{column}
  \end{columns}
\end{frame}

% ---------------------------------------------------------------------
\begin{frame}{Kesalahan Sintaksis dan Pemulihan}
  \begin{columns}[T,onlytextwidth]
    \begin{column}{0.46\textwidth}
      \textbf{Jenis kesalahan}
      \begin{itemize}
        \item titik koma hilang;
        \item kurung tidak seimbang;
        \item kata kunci salah;
        \item ekspresi tidak lengkap;
        \item deklarasi fungsi tanpa nama.
      \end{itemize}
    \end{column}
    \begin{column}{0.50\textwidth}
      \textbf{Strategi pemulihan}
      \begin{enumerate}
        \item \eng{panic mode}: lewati token sampai titik sinkronisasi;
        \item \eng{phrase-level recovery}: sisipkan token lokal;
        \item \eng{error production}: aturan khusus untuk kesalahan
              umum.
      \end{enumerate}
    \end{column}
  \end{columns}

  \vspace{4pt}
  \begin{block}{Mode panik pada predictive parser}
    \footnotesize
    \(H\) = token pembatas (titik koma, kurung) + token awal
    pernyataan. Jika \(M(A,a)=\text{error}\) dan \(a \in H\), lewati
    token sampai anggota \(H\). Jika \(a \in \mathrm{FOLLOW}(A)\),
    entri tabel menjadi \texttt{synch}.
  \end{block}
\end{frame}

% ---------------------------------------------------------------------
\begin{frame}{Contoh Pesan Kesalahan Sintaksis}
  \scriptsize
  \setlength{\tabcolsep}{4pt}
  \begin{tabular}{@{}p{3.0cm}p{10.4cm}@{}}
    \toprule
    \textbf{Masukan} & \textbf{Pesan} \\
    \midrule
    \texttt{int a = 1} &
      \texttt{baris 1, kolom 10: diharapkan ';' setelah deklarasi,
      ditemukan 'return'} \\
    \texttt{if (a} &
      \texttt{baris 1, kolom 6: diharapkan ')' setelah kondisi} \\
    \texttt{return ;} &
      \texttt{baris 1, kolom 8: ekspresi diharapkan setelah return} \\
    \texttt{int (} &
      \texttt{baris 1, kolom 5: nama fungsi diharapkan} \\
    \bottomrule
  \end{tabular}

  \vspace{4pt}
  {\scriptsize\color{ubbblue!75}Jika masukan terlalu rusak, parser
  berhenti setelah sejumlah kesalahan dan melaporkan ringkasan.}
\end{frame}

% ---------------------------------------------------------------------
\begin{frame}{Kegiatan Praktikum: M2}
  \begin{block}{Langkah kerja}
    \footnotesize
    \begin{enumerate}
      \setlength{\itemsep}{1pt}
      \item Hitung FIRST dan FOLLOW untuk grammar hasil Pertemuan 7.
      \item Bangun tabel LL(1) dan periksa konflik.
      \item Implementasikan fungsi parser per nonterminal.
      \item Bangun AST saat mengurangi aturan.
      \item Tambahkan sinkronisasi dan pesan kesalahan.
      \item Uji dengan program valid dan tidak valid.
    \end{enumerate}
  \end{block}

  \vspace{2pt}
  {\footnotesize\color{ubbblue!80}Kumpulkan juga artefak 3.1, 3.2, dan
  3.4 dari Pertemuan 6--7 karena Praktikum 2 menilai 3.1--3.4 dan
  4.1.}
\end{frame}

% ---------------------------------------------------------------------
\begin{frame}{Rubrik Praktikum 2}
  \small
  \setlength{\tabcolsep}{4pt}
  \begin{tabular}{@{}p{3.4cm}p{8.6cm}p{1.8cm}@{}}
    \toprule
    \textbf{Aspek} & \textbf{Indikator} & \textbf{Bobot} \\
    \midrule
    Kebenaran & Parser menerima program valid dan menolak yang tidak
      valid & 40\% \\
    Kualitas kode & Struktur fungsi, penamaan, komentar, tanpa
      peringatan & 25\% \\
    Kelengkapan & Ekspresi, pernyataan, AST, dan diagnostik & 20\% \\
    Pengujian & Kasus normal, batas, dan kesalahan terdokumentasi &
      15\% \\
    \bottomrule
  \end{tabular}

  \vspace{5pt}
  {\footnotesize\color{ubbblue!80}Rubrik ini sama dengan rubrik artefak
  praktikum pada bagian 10 RPS.}
\end{frame}

% ---------------------------------------------------------------------
\begin{frame}{Tugas dan Artefak}
  \begin{block}{Artefak Sub-CPMK 3.3 dan 4.1}
    Tabel FIRST/FOLLOW, keputusan LL(1), \textit{recursive descent
    parser}, dan \textit{predictive parser} LL(1) dalam bahasa C,
    disertai keluaran validasi.
  \end{block}

  \small
  \begin{itemize}
    \item Tabel FIRST/FOLLOW lengkap dan keputusan LL(1) beralasan.
    \item Parser C lulus minimal 20 kasus uji valid dan tidak valid.
    \item AST terbentuk untuk program valid.
    \item Diagnostik menyebut lokasi dan token yang diharapkan.
    \item Artefak 3.1, 3.2, dan 3.4 dari Pertemuan 6--7 ikut
          dikumpulkan.
    \item Laporan singkat diserahkan.
  \end{itemize}

  \vspace{3pt}
  {\footnotesize\color{ubbblue!75}Praktikum 2 berbobot 12\% dan
  dilaporkan setelah Milestone M2 selesai.}
\end{frame}

% ---------------------------------------------------------------------
\begin{frame}{Rangkuman dan Jembatan}
  \begin{itemize}
    \item Recursive descent memakai satu fungsi per nonterminal dan
          memilih aturan berdasarkan token saat ini.
    \item FIRST menentukan terminal awal suatu string; FOLLOW
          menentukan terminal setelah nonterminal.
    \item Grammar LL(1) jika tabel parsing tidak memiliki sel ganda.
    \item Predictive parser memakai stack eksplisit dan tabel LL(1)
          tanpa mundur.
    \item Pemulihan kesalahan memakai \eng{panic mode},
          \eng{phrase-level recovery}, atau \eng{error production}.
  \end{itemize}

  \vspace{4pt}
  \begin{exampleblock}{Pertemuan berikutnya}
    Pertemuan 10 beralih ke bottom-up parsing: \eng{shift-reduce} dan
    operator precedence. Pendekatan ini menangani \eng{left recursion}
    secara alami.
  \end{exampleblock}
\end{frame}

\end{document}
